
\begin{proposition}[Finite positive transport algebra]
\label{prop:finite-positive-transport-algebra}
Let
\[
C:=\frac12\bigl(U_6+U_6^\dagger\bigr),
\qquad
a:=y^2+\delta_{\mathrm{ph}}^2+\varepsilon^2,
\qquad
b:=y\delta_{\mathrm{ph}}.
\]
Then
\[
D_y^\dagger D_y+\varepsilon^2=a\,\mathbf 1-2b\,C.
\]
The cosine operator $C$ has spectrum
\[
\mathrm{Spec}(C)=\left\{1,\frac12,-\frac12,-1\right\}
\]
and satisfies
\begin{equation}
\label{eq:c6-cosine-polynomial}
4C^4-5C^2+\mathbf 1=0.
\end{equation}
Consequently the regularized positive transport kernel reduces to the cubic polynomial
\begin{equation}
\label{eq:finite-positive-transport-kernel}
\mathcal Y_y^{(\varepsilon)}
=
\frac{
a^3-5ab^2
+2b(a^2-5b^2)C
+4ab^2C^2
+8b^3C^3
}{
(a^2-b^2)(a^2-4b^2)
}.
\end{equation}
\end{proposition}

\begin{proof}
Because $U_6$ is unitary,
\[
D_y^\dagger D_y
=
\bigl(y\,\mathbf 1-\delta_{\mathrm{ph}}U_6^\dagger\bigr)
\bigl(y\,\mathbf 1-\delta_{\mathrm{ph}}U_6\bigr)
=
\bigl(y^2+\delta_{\mathrm{ph}}^2\bigr)\mathbf 1
-y\delta_{\mathrm{ph}}(U_6+U_6^\dagger),
\]
which gives the first identity after adding $\varepsilon^2$. The eigenvalues of $C$ are $\cos(n\pi/3)$ for $n=0,\dots,5$, namely $1,\frac12,-\frac12,-1$. Hence $C$ obeys the minimal polynomial \Cref{eq:c6-cosine-polynomial}. The inverse of $a-2bc$ on this four-point spectrum is therefore represented by a unique cubic interpolation polynomial in $c$, and evaluating that interpolation yields \Cref{eq:finite-positive-transport-kernel}.
\end{proof}

\begin{remark}[Carrier rank echo in the transport polynomial]\label{rem:transport-rank-echo}
On the minimal carrier branch $\gcar=5$, the hexagonal cosine polynomial of \Cref{eq:c6-cosine-polynomial} can be rewritten as
\[
(\gcar-1)C^4-\gcar C^2+\mathbf 1=0.
\]
Thus the transport closure polynomial carries the same pair $(\gcar-1,\gcar)$ that governs the hard carrier packet. Equivalently,
\[
4C^4-5C^2+\mathbf 1=(4C^2-\mathbf 1)(C^2-\mathbf 1)=(2C-\mathbf 1)(2C+\mathbf 1)(C-\mathbf 1)(C+\mathbf 1).
\]
So the distinguished value $C=1/2$ reappears on the transport side, mirroring the positive carrier root $Y=1/2$ of the factorized carrier polynomial in Section~3.3. At present this should be read as an exact compression identity on the retained hexagonal branch, not as a derivation of the transport sector from the carrier theorem.
\end{remark}

Absolute masses are attached to admissible path lengths,
\begin{equation}
\label{eq:path-length-yukawa}
y_{f,j}^{\mathrm{UV}}=\lambdaY^{L_{f,j}}\Lambda_{f,j},
\qquad
\lambdaY=\sqrt{\phiz(1-\phiz)}.
\end{equation}
The stronger matrix-level statement, used as the main transport holonomy grammar below, is that
\begin{equation}
\label{eq:yukawa-holonomy-matrix}
(Y_f)_{ij}
=
\lambdaY^{L^L_{f,i}+L^R_{f,j}}
\left\langle e_i,\mathrm{Hol}_{\Gamma_{ij}}(\nabla_F)e_j\right\rangle.
\end{equation}
\begin{proposition}[Exact family-holonomy phase]
\label{prop:holonomy-triangle-phase}
\label{prop:exact-family-holonomy-phase}
Define the ordered closed family holonomy by
\[
\Pi_\triangle
:=
\mathrm{Hol}_{\Gamma_{12}}(\nabla_F^\star)\,
\mathrm{Hol}_{\Gamma_{23}}(\nabla_F^\star)\,
\mathrm{Hol}_{\Gamma_{31}}(\nabla_F^\star),
\qquad
\delta_{\mathrm{CKM}}
:=
\arg\tr\Pi_\triangle.
\]
This exact holonomy phase is the benchmark quantity on the canonical admissible branch. On the
small transport-area branch one has
\begin{equation}
\label{eq:ckm-holonomy-phase}
\delta_{\mathrm{CKM}}
=
\frac{\pi}{3}+3\lambda_C^2+O(\lambda_C^4).
\end{equation}
\end{proposition}

\begin{proof}[Proof sketch]
The family holonomy is already encoded in \Cref{eq:yukawa-holonomy-matrix}. Closing the
minimal triangle isolates the gauge-invariant phase information carried by the family
connection. The $\mathbb Z_6$ transport grammar fixes the base phase at $\pi/3$, while the
first nontrivial nonabelian area correction is quadratic in the Cabibbo seed. Keeping only the
leading admissible area term gives the displayed asymptotic expansion of the exact phase.
\end{proof}

\Cref{eq:path-length-yukawa} is the diagonal or singular-value reduction of \Cref{eq:yukawa-holonomy-matrix}, not the conceptual starting point of the transport sector.
After the geometric electroweak closure fixes $\vgeo$, the intrinsic compiler mass rows are emitted through
\begin{equation}
\label{eq:mass-bridge}
\hat m_{f,j}:=\frac{y_{f,j}\vgeo}{\sqrt2}.
\end{equation}
These hatted quantities are the intrinsic compiler masses on the retained branch. The observation-surface masses are obtained only after the declared comparison map:
\[
m_{f,j}^{\mathrm{obs}}:=\mathcal R_{\SM}(\hat m_{f,j}).
\]
The role of the transport matrix elements $\Lambda_{f,j}$ must be kept explicit. They are bounded $O(1)$ admissible amplitudes, not unrestricted fit knobs. The RG bookkeeping behind direct $\vgeo$ insertions is part of the declared comparison map discussed later in \Cref{sec:rg-map-compiler-outputs}.

\begin{remark}[Singlet-projected cusp lift to the transport hexagon]
The transport cusps do not come from the raw two-point carrier spectrum alone. The carrier itself has
\[
\mathrm{Spec}(Y|_E)=\left\{-\frac13,\frac12\right\}.
\]
After exterior lift to the one-family packet,
\[
\mathrm{Spec}(Y|_{S^+})
=
\left\{0,-\frac23,\frac16,1,-\frac12,\frac13\right\},
\]
with multiplicities displayed in \Cref{tab:one-family}. Restricting to the singlet channels and taking absolute values yields
\[
\{|Y(e^c)|,\ |Y(u^c)|,\ |Y(d^c)|\}
=
\left\{1,\frac23,\frac13\right\},
\]
which are exactly the admissible cusp values of the transport hexagon. In this sense the transport sector reads the singlet-projected, family-lifted hypercharge spectrum through the resolvent
\[
D_y=y\mathbf 1-\delta_{ph}U_6
\]
on $C_6$. The transport grammar therefore lives on the same rational grid as the carrier packet, but only after exterior lift and singlet projection.
\end{remark}

On the canonical admissible branch selected numerically, the retained diagonal word-length packets are
\begin{equation}
\label{eq:minimal-word-lengths}
\begin{gathered}
L_u^{\mathrm{diag}}=(L_u,L_c,L_t)=(7,3,0),\\
L_d^{\mathrm{diag}}=(L_d,L_s,L_b)=(7,5,2),\\
L_e^{\mathrm{diag}}=(L_e,L_\mu,L_\tau)=(8,5,3).
\end{gathered}
\end{equation}
These are not free parameters; they are the canonical retained diagonal word-length packets on the numerically selected admissible branch, with transport matrix elements in \Cref{tab:yukawa-residuals} remaining in the bounded $O(1)$ range. Once the hierarchy ordering is fixed, the admissible graph distance is determined up to full $6$-step windings by $\lambdaY\approx\num{0.2244}$.

\begin{remark}[Residue classes and full windings on the hexagon]
The exact resolvent formula \Cref{eq:finite-hexagonal-resolvent} makes the combinatorics behind \Cref{eq:minimal-word-lengths} transparent: every path length splits into a residue class on the base hexagon plus an optional full winding by $6$. Accordingly,
\[
L_t=0,\qquad
L_b=2,\qquad
L_\tau=L_c=3,\qquad
L_s=L_\mu=5
\]
live on the first hexagon, while
\[
\;L_u=L_d=1+6,
\qquad
L_e=2+6
\]
carry one additional full turn. Since
\[
\lambda_Y^6=\bigl(\phiz(1-\phiz)\bigr)^3\approx \num{1.28e-4},
\]
each extra winding produces the severe suppression needed for the first generation. This is not by itself a full mass theorem, but it is an exact algebraic reading of why the retained word-length pattern is so uneven.
\end{remark}

\begin{remark}[Universal $C_6$ transport triangle]
Modulo full windings, the three diagonal packets define the residue sets
\[
\{L_t,L_c,L_u\}\equiv\{0,3,1\},
\]
\[
\{L_b,L_s,L_d\}\equiv\{2,5,1\}=2-\{0,1,3\},
\qquad
\{L_\tau,L_\mu,L_e\}\equiv\{3,5,2\}=2+\{0,1,3\}
\pmod 6.
\]
Hence the up, down, and charged-lepton sectors lie on the same $D_6$ orbit of the base three-point subset $\{0,1,3\}\subset C_6$. In each case the cyclic distance multiset is exactly
\[
\{1,2,3\}.
\]
So the exact structure here is a three-point circular distance system with pairwise distinct cyclic distances, not a classical perfect difference set or a linear Golomb ruler.
Moreover, componentwise division by the full winding length gives
\[
\left\lfloor\frac{L_u^{\mathrm{diag}}}{6}\right\rfloor
=
\left\lfloor\frac{L_d^{\mathrm{diag}}}{6}\right\rfloor
=
\left\lfloor\frac{L_e^{\mathrm{diag}}}{6}\right\rfloor
=(1,0,0),
\]
so the first-generation direction carries exactly one extra full winding in every sector. This
shows that the closed admissible branch preserves one universal $C_6$ transport triangle together
with a common first-generation winding lock.
\end{remark}

\begin{corollary}[Diagonal transport packets on the selected branch]
\label{prin:minimal-admissible-transport-packets}
On the canonical admissible branch selected numerically, the retained diagonal transport packets
are
\[
L_u^{\mathrm{diag}}=(7,3,0),
\qquad
L_d^{\mathrm{diag}}=(7,5,2),
\qquad
L_e^{\mathrm{diag}}=(8,5,3).
\]
\end{corollary}

\begin{corollary}[Diagonal transport sum rule on the canonical branch]
\label{cor:transport-length-sum-rule}
Write
\[
L_f^{\mathrm{diag}}=(L_{f,1}^{\mathrm{diag}},L_{f,2}^{\mathrm{diag}},L_{f,3}^{\mathrm{diag}})
\qquad
(f\in\{u,d,e\}).
\]
Then on the canonical branch
\[
\sum_{f\in\{u,d,e\}}\sum_{j=1}^{3}L_{f,j}^{\mathrm{diag}}
=
\Oadm\gamma
=
10\,\bone-1
=
40.
\]
\end{corollary}

\begin{proof}
By \Cref{prin:minimal-admissible-transport-packets},
\[
\sum_{j=1}^3L_{u,j}^{\mathrm{diag}}=7+3+0=10,
\qquad
\sum_{j=1}^3L_{d,j}^{\mathrm{diag}}=7+5+2=14,
\]
\[
\sum_{j=1}^3L_{e,j}^{\mathrm{diag}}=8+5+3=16.
\]
Hence
\[
\sum_{f,j}L_{f,j}^{\mathrm{diag}}=10+14+16=40.
\]
On the other hand, \Cref{eq:41-chain-main} gives
\[
\Oadm\gamma=10\,\bone-1=40.
\]
Therefore the two quantities agree exactly.
\end{proof}

\begin{remark}[Transport sum rule on the closed branch]
The equality of \Cref{cor:transport-length-sum-rule} is now read as an exact cross-layer sum
rule on the closed transport branch rather than as a still-open dynamical upgrade path.
\end{remark}

The transport-side status is summarized in \Cref{tab:transport-dictionary}.

\begin{table}[t]
\centering
\scriptsize
\resizebox{\linewidth}{!}{%
\begin{tabular}{@{}lllll@{}}
\toprule
Sector & Admissible length data & Transport matrix element & Role & Status \\
\midrule
charged leptons & $L_e^{\mathrm{diag}}=(8,5,3)$ & $\Lambda_e,\Lambda_\mu,\Lambda_\tau$ & ladder and hierarchy & conditional closure \\
up quarks & $L_u^{\mathrm{diag}}=(7,3,0)$ & $\Lambda_u,\Lambda_c,\Lambda_t$ & up-type ratio ladder & conditional closure \\
down quarks & $L_d^{\mathrm{diag}}=(7,5,2)$ & $\Lambda_d,\Lambda_s,\Lambda_b$ & down-type ratio ladder & conditional closure \\
quark mixing & $L_{ij}^{q}$ & $\Lambda_{ij}^{q}$ & CKM transport entries & matched intrinsic observation \\
neutrino mixing & $L_{ij}^{\nu}$ & $\Lambda_{ij}^{\nu}$ & PMNS transport entries & matched intrinsic observation \\
hadron basis & center-neutral word lengths & block coefficients & meson and baryon transport basis & singlet closure on $\Hhad$ \\
strong CP & determinant spectrum of $\mathcal Y_f$ & none on the positive branch & phase suppression & conditional closure \\
\bottomrule
\end{tabular}}
\caption{Transport data on the selected branch, with status labels kept separate between conditional closure statements and matched observation-facing rows.}
\label{tab:transport-dictionary}
\end{table}

\subsection{CKM and PMNS sharpening}

The leading CKM magnitudes on the transport branch are
\begin{equation}
\label{eq:ckm-magnitudes}
s_{12}=\lambda_C,
\qquad
s_{23}=\frac{\phiz}{1+\lambda_C},
\qquad
s_{13}=\frac{\lambda_C^3}{3}.
\end{equation}
As a small-area transport asymptotic, \Cref{prop:holonomy-triangle-phase} gives
\begin{equation}
\label{eq:ckm-phase}
\delta_{\mathrm{CKM}}=\frac{\pi}{3}+3\lambda_C^2+O(\lambda_C^4).
\end{equation}
The benchmark row used later uses the exact holonomy phase
\[
\delta_{\mathrm{CKM}}=\arg\tr\Pi_\triangle(\nabla_F^\star),
\]
while \Cref{eq:ckm-phase} is retained only as its small-area asymptotic expansion. Thus
$\pi/3$ remains the base monodromy phase while $3\lambda_C^2$ is the leading
family-transport area correction.

\begin{theorem}[Neutrino closure from the same transport branch]
The Majorana neutrino sector is generated by the same admissible transport grammar together
with the seam-even Majorana pairing on the family bundle. Hence the light-neutrino mass
matrix, PMNS matrix, and mass sum are intrinsic outputs of the same admissible branch and are
not imported from an external oscillation-inversion package.
\end{theorem}

\begin{remark}[Carrier-complement reading of the neutrino hierarchy]
On the canonical admissible branch the leading neutrino hierarchy ratio still admits the
compact carrier reading
\[
\frac{m_2}{m_3}\sim \pi\phiz
=
\pi\!\left(\frac{1-\gamma}{\pi}+\frac{3}{256\pi^4}\right)
=
(1-\gamma)+\frac{3}{256\pi^3},
\]
so the carrier complement $1-\gamma=1/6$ remains the dominant organizing term, corrected by a
small seam contribution.
\end{remark}

\subsection{Confinement surrogate as admissible singlet selection}

\begin{proposition}[Admissible color words]
For each carrier word $w$, let $z(w)\in\{1,\omega,\omega^2\}$ be its color-center charge and let $\varepsilon(w)\in\{\pm1\}$ be its seam parity. Physical hadronic words satisfy
\[
z(w)=1,
\qquad
\varepsilon(w)=+1.
\]
\end{proposition}

This is a superselection statement on the retained branch, not a derivation of dynamical confinement from the microscopic Yang--Mills flow. The corresponding hadronic Hilbert space will be written explicitly in \Cref{sec:closure-theorems}. A schematic version is shown in \Cref{fig:admissible-words}.

\begin{figure}[t]
\centering
\begin{tikzpicture}[>=Latex, node distance=1.1cm and 1.3cm]
\node[draw=red!70!black, fill=red!4, rounded corners, minimum width=2.1cm, minimum height=1.0cm, align=center] (q) {free quark\\$z(w)\neq 1$};
\node[draw=green!50!black, fill=green!4, rounded corners, minimum width=2.1cm, minimum height=1.0cm, align=center, right=of q] (m) {meson\\$z(w)=1$};
\node[draw=green!50!black, fill=green!4, rounded corners, minimum width=2.1cm, minimum height=1.0cm, align=center, right=of m] (b) {baryon\\$z(w)=1$};
\draw[red!70!black, very thick] ($(q.north west)+(0.12,-0.12)$) -- ($(q.south east)+(-0.12,0.12)$);
\draw[red!70!black, very thick] ($(q.south west)+(0.12,0.12)$) -- ($(q.north east)+(-0.12,-0.12)$);
\node[below=0.2cm of q] {\small excluded};
\node[below=0.2cm of m] {\small admissible};
\node[below=0.2cm of b] {\small admissible};
\end{tikzpicture}
\caption{Confinement as admissible word selection.}
\label{fig:admissible-words}
\end{figure}

The transport-sector string tension is written as
\[
\sigma=\cthree^2\chi_{\QCD}^2,
\]
so the relative Wilson loop obeys
\[
\langle W(C)\rangle_{\mathrm{rel}}\sim e^{-\sigma A(C)}.
\]
The sector-dependent scale $\chi_{\QCD}$ is the minimal notation needed to avoid confusing the hadronic transport channel with the gravitational $\chigeozero$ background; the two channels are linked by the common underlying primitive response $\chiseed$ but not by a naive equality of scales.

